\begin{ThreePartTable}

\begin{TableNotes}
    \item[*] 所有内部存储器均接受权限管理。只有获取到访问内部存储器的访问权限，才可以执行正常的访问操作，CPU 访问内部存储器时才可以被响应。关于权限管理的更多信息，请参考章节 \ref{mod:permctrl} \textit{\nameref{mod:permctrl}}。

\end{TableNotes}

\begin{longtable}[c]{|p{3.5cm}|c|c|r|c| } 
    \caption{地址映射}
    \label{tab:sysmem-memory-cn}
       \\\hline
        \rowcolor{lightgray}
       & \multicolumn{2}{c|}{\textbf{边界地址}} & & \\\cline{2-3}
       
       \rowcolor{lightgray}
       \multirow{-2}*{\textbf{总线类型}} & \textbf{低位地址} & \textbf{高位地址} & \multirow{-2}*{\textbf{容量}} & \multirow{-2}*{\textbf{目标}} \\\hline
    \endfirsthead
    
       \multicolumn{5}{c}{\bfseries \tablename\ \thetable{} -- 接上页}\\\hline
       
       \rowcolor{lightgray}
       & \multicolumn{2}{c|}{\textbf{边界地址}} & & \\\cline{2-3}
       
       \rowcolor{lightgray}
       \multirow{-2}*{\textbf{总线类型}} & \textbf{低位地址} & \textbf{高位地址} & \multirow{-2}*{\textbf{容量}} & \multirow{-2}*{\textbf{目标}} \\\hline
    \endhead
    
        \multicolumn{5}{r}{\textbf{见下页}} \\
    \endfoot
    
    \insertTableNotes
    \endlastfoot
        \rowcolor{gray!25}  & 0x0000\_0000 & 0x1FFF\_FFFF &         & 保留      \\\hline
        {数据总线/指令总线} & 0x2000\_0000 & 0x2FFF\_FFFF & 256 MB  & CPU Sub-system \\\hline

        \hline
        \rowcolor{gray!25}  & 0x3000\_0000 & 0x3FFF\_FFFF &        & 保留      \\\hline
        {数据总线/指令总线} & 0x4000\_0000 & 0x4001\_FFFF & 128 KB  & ROM\tnote{*}  \\\hline
        \hline

        \rowcolor{gray!25}  & 0x4002\_0000 & 0x407F\_FFFF &         & 保留      \\\hline
        {数据总线/指令总线} & 0x4080\_0000 & 0x4084\_FFFF & 320 KB  & HP SRAM\tnote{*} \\\hline

        \hline

        \rowcolor{gray!25}  & 0x4085\_0000 & 0x41FF\_FFFF &         & 保留      \\\hline
        {数据总线/指令总线} & 0x4200\_0000 & 0x42FF\_FFFF & 16 MB & 外部存储器  \\\hline
        \rowcolor{gray!25}  & 0x4300\_0000 & 0x4FFF\_FFFF &         & 保留      \\\hline
        {数据总线/指令总线} & 0x5000\_0000 & 0x5000\_0FFF & 4 KB  & LP SRAM\tnote{*} \\\hline 
        \rowcolor{gray!25}  & 0x5000\_1000 & 0x5FFF\_FFFF &         & 保留      \\\hline
        {数据总线/指令总线} & 0x6000\_0000 & 0x600C\_FFFF & 832 KB  & 外设 \\\hline
        \rowcolor{gray!25}  & 0x600D\_0000 & 0xFFFF\_FFFF &         & 保留      \\\hline
        
\end{longtable}
\end{ThreePartTable}


